\section*{Limitations}

In the following, we outline the limitations of our work to ensure transparency and inspire future research.
First, the chart domains we experimented with is limited to a few popular websites (e.g. Statista and Pew). This is due to the fact that existing academic chart-to-summary datasets only cover limited domains. However, to comprehensively evaluate the effectiveness of \metric{} and \pipeline{}, it is desirable to also evaluate our approaches in other chart domains such as infographics and scientific/financial charts.
Second, the \metric{} depends on a deplotter (image-to-text) model, specifically DePlot~\cite{liu-etal-2023-deplot}. DePlot has been trained on similar domains as the chart-to-summary datasets used in this work (e.g. Statista), and its performance may not generalize to other domains. The SciCap dataset evaluated in \Cref{tab:nlg-llm_results} provides some evidence of generalization. 
In future work, we plan to build out-of-domain evaluations to understand the impact of the deplotter's robustness better.
Third, we focused mostly on English chart summary in this work. We plan to also expand multilingual chart summary in future works and expanding our evaluation on the TaTa dataset \citep{gehrmann2022tata} as a test bed.

Another potential limitation is the use of \metric{} to evaluate \pipeline{}, although it is important to note that in order to mitigate the possibility for bias, we have cross checked with a variety of evaluation metrics and conducted ablation studies to analyze the individual components of our proposed solution. 
\pipeline{} has proven to also provide significant benefits when employing standard reference-based metrics. 
Comprehensive human studies have corroborated the usefulness of our results and we can confidently say that the improvement is consistent across the board.

We would also like to highlight the underlying risk of blindly trusting models to summarize content from an image accurately. Special care should be taken to verify outputs in accuracy-sensitive applications.

Despite its limitations, our work serves as an initial step in constructing reliable chart summarization evaluations and models. We hope future research can greatly benefit from this starting point.

\section*{Acknowledgments}
We would like to thank Massimo Nicosia,
Srini Narayanan, 
Yasemin Altun,
Averi Nowak,
Chenxi Pang,
Jonas Pfeiffer,
Mubashara Akhtar,
Parag Jain,
and the anonymous reviewers for their time, constructive feedback, useful
comments and suggestions about this work.
